% OVERLEAF: New Project > Blank > coller ce fichier dans main.tex > pdfLaTeX
% Offre: MARSO Robotics — Embedded Software Engineer (Paris, 45-60k)
\documentclass[10pt,a4paper]{article}
\usepackage[margin=1.2cm]{geometry}
\usepackage[T1]{fontenc}
\usepackage[utf8]{inputenc}
\usepackage{hyperref,enumitem,titlesec}
\hypersetup{hidelinks}
\pagestyle{empty}
\setlength{\parindent}{0pt}
\setlength{\parskip}{0pt}
\titleformat{\section}{\bfseries\large\uppercase}{}{0em}{}[\vspace{1pt}\titlerule\vspace{3pt}]
\titlespacing*{\section}{0pt}{6pt}{3pt}
\setlist[itemize]{leftmargin=*,itemsep=1pt,topsep=1pt,parsep=0pt}
\begin{document}

\begin{center}
{\LARGE\bfseries Ghislain THEUBO}\\[2pt]
{\large Ing\'enieur Logiciel Embarqu\'e --- C/C++ / Linux / Temps r\'eel}\\[3pt]
+33 6 82 86 62 79 $\cdot$
\href{mailto:bergerghislaintheubo@gmail.com}{bergerghislaintheubo@gmail.com} $\cdot$
\href{https://www.linkedin.com/in/berger-ghislain}{LinkedIn} $\cdot$
\href{https://github.com/Bergerghislain}{GitHub}\\[1pt]
Paris (pr\'esentiel) $\cdot$ Dispo imm\'ediate $\cdot$ Objectif salarial 48--55\,k€ $\cdot$ EN/FR
\end{center}

\section{Profil}
Ing\'enieur Bac+5 (Grenoble INP -- Ensimag) sp\'ecialis\'e logiciel, avec une premi\`ere exp\'erience concr\`ete en \textbf{d\'eveloppement embarqu\'e temps r\'eel} (INRIA : C/FreeRTOS, Python, optimisation de latence sur syst\`eme r\'eel). Solide pratique de \textbf{C/C++}, \textbf{Python}, \textbf{Linux}, et culture qualit\'e / tests en environnement industriel (Dassault). Motiv\'e pour rejoindre MARSO : concevoir, int\'egrer et optimiser le logiciel embarqu\'e qui relie perception / IA et hardware robotique en conditions terrain (entrep\^ots brownfield).

\section{Exp\'eriences professionnelles}
\textbf{INRIA Grenoble --- Humanlab (projet Froll'N'Roll)} \hfill Montbonnot-Saint-Martin\\
\textit{Stagiaire --- Logiciel embarqu\'e temps r\'eel (C/FreeRTOS, Python)} \hfill Mai 2024 -- Ao\^ut 2024
\begin{itemize}
  \item Conception et d\'eveloppement de modules \textbf{C / FreeRTOS} pour le contr\^ole interactif d'un syst\`eme physique (contraintes de latence et de fiabilit\'e).
  \item Traitement de \textbf{flux capteurs} et pipelines applicatifs ; outillage et scripts de test en \textbf{Python}.
  \item Refonte de protocoles / structures de donn\'ees : gain d'environ \textbf{40\,\%} de r\'eactivit\'e mesur\'ee.
  \item Validation continue avec utilisateurs finaux ; it\'eration ship $\rightarrow$ mesure $\rightarrow$ correction (esprit tests terrain).
  \item Collaboration pluridisciplinaire (recherche, hardware-adjacent, UX) --- interface bas niveau $\leftrightarrow$ couches applicatives.
\end{itemize}

\textbf{Gustave Roussy \& CentraleSup\'elec MICS} \hfill Paris / Saclay\\
\textit{Software Engineer (CDD)} \hfill Jan. 2026 -- Sept. 2026
\begin{itemize}
  \item Services \textbf{Python} en production, tests, \textbf{Docker} et CI (GitHub Actions).
  \item Validation rigoureuse avec op\'erateurs m\'etier ; collecte de feedback et correction d'anomalies.
  \item Collaboration avec \'equipes data / IA --- pont naturel avec le volet ``AI collaboration'' de MARSO.
\end{itemize}

\textbf{Dassault Syst\`emes --- 3DEXPERIENCE / 3DSwym R\&D} \hfill V\'elizy-Villacoublay\\
\textit{Stagiaire Software Engineer} \hfill Mars 2025 -- Sept. 2025
\begin{itemize}
  \item Exigence industrielle : optimisation de chemins critiques ($\sim$\textbf{35\,\%}), couverture de tests $\sim$\textbf{85\,\%}, TDD.
  \item Documentation, revues de code, livraison en \'equipe R\&D --- habitudes de robustesse pour syst\`emes complexes.
\end{itemize}

\textbf{LIG --- Laboratoire d'Informatique de Grenoble} \hfill Grenoble\\
\textit{Stagiaire infrastructure cloud / Linux} \hfill F\'evr. 2025 -- Mars 2025
\begin{itemize}
  \item Automatisation Linux, mesures de performance, environnements reproductibles --- base pour bring-up / debug Ubuntu embarqu\'e.
\end{itemize}

\textbf{Shift Technology --- Platform Software Engineering} \hfill Paris\\
\textit{Software Engineer (CDI)} \hfill Sept. 2026 -- Nov. 2026
\begin{itemize}
  \item Culture ownership en production ; communication transversale claire (utile en startup deeptech).
\end{itemize}

\section{Formation}
\textbf{Grenoble INP -- Ensimag} \hfill 2023 -- 2025\\
Dipl\^ome d'ing\'enieur --- Syst\`emes d'information / G\'enie logiciel\\[2pt]
\textbf{\'Ecole Nationale Sup\'erieure Polytechnique de Yaound\'e} \hfill 2019 -- 2023\\
Master G\'enie informatique --- Top 12\,\% de la promotion

\section{Comp\'etences techniques}
\textbf{Langages :} C, C++, Python, notions scripts / outillage de test\\
\textbf{Embarqu\'e :} FreeRTOS, Linux / Ubuntu, contraintes temps r\'eel, optimisation de latence\\
\textbf{Cible MARSO (mont\'ee en comp\'etence) :} ROS / ROS2, drivers bus (UART, SPI, I2C, CAN), capteurs (cam\'era, IMU, LiDAR), r\'epartition CPU/GPU\\
\textbf{M\'ethodes :} Git, CI, Docker, tests, qualit\'e pragmatique, am\'elioration continue

\section{Langues}
\textbf{Anglais} : bilingue $\cdot$ \textbf{Fran\c{c}ais} : courant $\cdot$ Esprit d'\'equipe, curiosit\'e tech, orientation terrain

\end{document}
